\documentclass{article}
\usepackage[utf8]{inputenc}
\usepackage[T1]{fontenc}
\usepackage{geometry}
\geometry{
    a4paper,
    left=2.5cm,
    right=2.5cm,
    top=2.5cm,
    bottom=2.5cm
}
\usepackage{graphicx}
\usepackage{float}
\usepackage{hyperref}
\usepackage[english]{babel}
\usepackage{amsmath}
\usepackage{amssymb}
\usepackage{array}
\usepackage{booktabs}
\usepackage{caption}
\usepackage{subcaption}
\usepackage{enumitem}
\usepackage{fancyhdr}
\usepackage{titlesec}
\usepackage{url}
\usepackage{textcomp}
\usepackage{xcolor}
\usepackage{listings}
\usepackage{multirow}
\usepackage{longtable}
\usepackage{tabularx}
\usepackage{times}

% Set up page style
\pagestyle{fancy}
\fancyhf{}
\rfoot{\thepage}
\renewcommand{\headrulewidth}{0pt}

% Title formatting
\titleformat{\chapter}[display]{\normalfont\bfseries}{\chaptertitle\thechapter}{0pt}{\Huge}
\titlespacing*{\chapter}{0pt}{50pt}{40pt}

% Document info
\title{Template Document}
\author{Academic Support Agent}
\date{\today}

\begin{document}
\maketitle

\pagenumbering{roman} % Se utilizan cifras romanas en la numeraci n de las p ginas previas al cuerpo del trabajo
	
%----------
%	PORTADA
%----------	
\begin{titlepage}
	\begin{sffamily}
	\color{azulUC3M}
	\begin{center}
		\begin{figure}[H] %incluimos el logotipo de la Universidad
			\makebox[\textwidth][c]{\includegraphics[width=16cm]{logo_UC3M.png}}
		\end{figure}
		\vspace{2.5cm}
		\begin{Large}
			Grado Universitario...\\			
			 2020-2021\\ %Indica el curso acad mico
			\vspace{2cm}		
			\textsl{Trabajo Fin de Grado}
			\bigskip
			
		\end{Large}
		 	{\Huge ``T tulo del trabajo''}\\
		 	\vspace*{0.5cm}
	 		\rule{10.5cm}{0.1mm}\\
			\vspace*{0.9cm}
			{\LARGE Nombre Apellido1 Apellido2}\\ 
			\vspace*{1cm}
		\begin{Large}
			Tutor/es\\
			Nombre Apellido1 Apellido2\\
			Nombre Apellido1 Apellido2\\
			Lugar y fecha de presentaci n prevista\\
		\end{Large}
	\end{center}
	\vfill
	\color{black}
	% SI NUESTRO TRABAJO SE VA A PUBLICAR CON UNA LICENCIA CREATIVE COMMONS, INCLUIR ESTAS L NEAS. ES LA OPCI N RECOMENDADA.
	\includegraphics[width=4.2cm]{creativecommons.png}\\ %incluimos el logotipo de Creative Commons
	Esta obra se encuentra sujeta a la licencia Creative Commons \textbf{Reconocimiento - No Comercial - Sin Obra Derivada}
	\end{sffamily}
\end{titlepage}

\newpage %p gina en blanco o de cortes a
\thispagestyle{empty}
\mbox{}

%----------
%	RESUMEN Y PALABRAS CLAVE
%----------	
\renewcommand\abstractname{\large\bfseries\filcenter\uppercase{Resumen}}
\begin{abstract}
\thispagestyle{plain}
\setcounter{page}{3}
	
	% ESCRIBIR EL RESUMEN AQU 
	
	\textbf{Palabras clave:}
	% Escribir las palabras clave aqu 
	
	\vfill
\end{abstract}
	\newpage % p gina en blanco o de cortes a
	\thispagestyle{empty}
	\mbox{}


%----------
%	DEDICATORIA
%----------	
\chapter*{Dedicatoria} % \chapter* evita que aparezca en el  ndice

\setcounter{page}{5}
	
	% ESCRIBIR LA DEDICATORIA AQU 	
		
	\vfill
	
	\newpage % p gina en blanco o de cortes a
	\thispagestyle{empty}
	\mbox{}
	

%----------
%	 NDICES
%----------	

%--
%  ndice general
%-

\thispagestyle{fancy}

\newpage % p gina en blanco o de cortes a
\thispagestyle{empty}
\mbox{}

%--
%  ndice de figuras. Si no se incluyen, comenta las l neas siguientes
%-

\thispagestyle{fancy}

\newpage % p gina en blanco o de cortes a
\thispagestyle{empty}
\mbox{}

%--
%  ndice de tablas. Si no se incluyen, comenta las l neas siguientes
%-

\thispagestyle{fancy}

\newpage % p gina en blanco o de cortes a
\thispagestyle{empty}
\mbox{}


%----------
%	MEMORIA
%----------	

\pagenumbering{arabic} % numeraci n con n meros ar bigos para el resto de la memoria.	

\chapter{Introducci n}

	% COMENZAR A ESCRIBIR la MEMORIA
	
	% IMPORTANTE: en LaTeX hay una serie de caracteres especiales, que son: # $ % & \ ^ _ { } ~. Si aparecen en el texto, tendr s un error al compilar. La mayor a se pueden escapar escribiendo \ delante. Para \ utiliza \textbackslash ; para ^ \textasciitilde y para ~ \textasciicircum.

    % COMO incluir una FIGURA siguiendo las recomendaciones de la Gu a: Alineaci n del t tulo: izquierda, en la parte superior de la tabla; fuente a 10pt (el resto de la memoria est  a 12); Numeraci n de la figura en negrita, t tulo de la figura en nueva l nea en cursiva; Propiedad intelectual: Se debe indicar la fuente de origen de la informaci n en la parte inferior de la figura, a continuaci n del t tulo.
    
    % EJEMPLO DE INCLUSI N de una figura:
    % \begin{figure}[H]
    % 	\ffigbox[\FBwidth] {
    % 	\caption[Nombre que aparecer  en el  ndice]{Nombre que aparecer  sobre la figura}
    % 	\floatfoot{ d nde hemos obtenido la imagen?}% P.e.: \floatfoot{Fuente: \textcite[p. 35]{id-referencia}}
    %     }
    % 	{\includegraphics[scale=0.6]{imagenes/creativecommons.png}}
    % \end{figure}
    
    % COMO incluir una TABLA siguiendo las recomendaciones de la Gu a: Alineaci n del t tulo: izquierda, en la parte superior de la tabla; fuente a 10pt (el resto de la memoria est  a 12); Numeraci n de la figura en negrita, t tulo de la figura en nueva l nea en cursiva; Propiedad intelectual: Se debe indicar la fuente de origen de la informaci n en la parte inferior de la tabla, a continuaci n del t tulo.
    
    % EJEMPLO DE INCLUSI N de una tabla:
% \begin{table}[H]
% 	\ttabbox[\FBwidth]
% 	{\caption{Lorem ipsum}
% 	\floatfoot{Fuente: BOE}}
% 	{\begin{tabular}{|c|P{1.5cm}|c|P{1.5cm}|P{2cm}|c|P{1.5cm}|P{2cm}|}
% 		\hline
% 		\multicolumn{2}{|c|}{\textbf{I}} & \multicolumn{2}{c|}{\textbf{II}} & \multicolumn{3}{c|}{\textbf{III}} & \textbf{IV} \\
% 		\hline
% 		x & y & x & y & x & y & x & y \\
% 		\hline
% 		10.0 & 8.04 & 10.0 & 9.14 & 10.0 & 7.46 & 8.0 & 6.58 \\
% 		\hline
% 		8.0 & 6.95 & 8.0 & 8.14 & 8.0 & 6.77 & 8.0 & 5.76 \\
% 		\hline
% 		13.0 & 7.58 & 13.0 & 8.74 & 13.0 & 12.74 & 8.0 & 7.71 \\
% 		\hline
% 		9.0 & 8.81 & 9.0 & 8.77 & 9.0 & 7.11 & 8.0 & 8.84 \\
% 		\hline
% 		11.0 & 8.33 & 11.0 & 9.26 & 11.0 & 7.81 & 8.0 & 8.47 \\
% 		\hline
% 		14.0 & 9.96 & 14.0 & 8.10 & 14.0 & 8.84 & 8.0 & 7.04 \\
% 		\hline
% 		6.0 & 7.24 & 6.0 & 6.13 & 6.0 & 6.08 & 8.0 & 5.25 \\
% 		\hline
% 		4.0 & 4.26 & 4.0 & 3.10 & 4.0 & 5.39 & 19.0 & 12.50 \\
% 		\hline
% 		12.0 & 10.84 & 12.0 & 9.13 & 12.0 & 8.15 & 8.0 & 5.56 \\
% 		\hline
% 		7.0 & 4.82 & 7.0 & 7.26 & 7.0 & 6.42 & 8.0 & 7.91 \\
% 		\hline
% 		5.0 & 5.68 & 5.0 & 4.74 & 5.0 & 5.73 & 8.0 & 6.89 \\
% 		\hline
% 	\end{tabular}}
% \end{table}


%----------
%	BIBLIOGRAF A
%----------	

%\nocite{*} % Si quieres que aparezcan en la bibliograf a todos los documentos que la componen (tambi n los que no est n citados en el texto) descomenta est  l nea


\addcontentsline{toc}{chapter}{Bibliograf a}
\setquotestyle[english]{british} % Cambiamos el tipo de cita porque en el estilo IEEE se usan las comillas inglesas.
\printbibliography



%----------
%	ANEXOS
%----------	

% Si tu trabajo incluye anexos, puedes descomentar las siguientes l neas
%\chapter* {Anexo x}
%\pagenumbering{gobble} % Las p ginas de los anexos no se numeran
\end{document}